\documentclass[12px]{article}

\title{Lezione 7 Metodi}
\date{2026-03-20}
\author{Federico De Sisti}

\input{../../setup.tex}

\begin{document}
	\maketitle
	\newpage
	\section{Convergenza}
	\subsection{Riassuntino}
	\textbf{Consistenza:} Lo schema approssima l'ODE? Se $Z(h) \rightarrow 0$ (l'errore tende a 0)\\
	o alternativamente $\Phi(t_n,u_n) \xrightarrow{h \rightarrow 0}f(t_n,u_n)$.\\
	Si sostituisce la soluzione esatta $y$ di ODE nello schema e usando sviluppi di Taylor (assumendo che la soluzione sia sufficientemente regolare)\\
	si mostra che  $\tau(h) \rightarrow 0$\\
	Se $y$ è sufficientemente regolare $\tau(h) = O(h^q)$ ($q$ è ordine di consistenza)\\[10px]
	\textbf{Zero-stabilità:} Lo schema è poco sensibile alle perturbazioni.\\
$|u_n-z_n|\leq C|u_0-z_0| = C\varepsilon$   (differenza  tra soluzione perturbata e soluzione non perturbata, controllata dalla differenza iniziale)\\
quindi se $C\varepsilon < tol \Rightarrow  \varepsilon < \frac{tol}C$ \\
Se $C$ diventa molto grande sono costretto ad un  $\varepsilon$ molto piccolo.\\
Se $\Phi$ è Lipschitziana con costante $\Lambda \Rightarrow $ lo schema (il metodo) è zero-stabile.\\
\subsection{Convergenza}\\
$u_n \xrightarrow{ h \rightarrow 0 } y_n$
\begin{defi}[Convergenza]
	Uno schema si dice convergente se $\exists C = C(h)$ tale che 
	\[
		|u_n - y_n|\leq C(h) \ \ \text{ e } \ \ C(h) \xrightarrow{h \rightarrow 0} 0\ \ \ \ 0\leq n\leq N
	.\] 
Se $y$ è sufficientemente regolare  e $C(h) = O(h^r)$ allora il metodo si dice convergente con ordine $r$. Chiameremo $r$ ordine di convergenza.
\end{defi}
\begin{teo}[di equivalenza di Lax-Richtmyer (Dahlquist)]
	Se lo schema numerico è consistente e $|y_0-u_0| \xrightarrow{h \rightarrow 0} 0$ allora è convergente.\\
	Se inoltre $|y_0 - u_0| = O(h^s)$ e lo schema ha ordine di consistenza $O(h^q)$ allora è convergente con ordine  $O(h^{\min\{s,q\}})$
\end{teo}
\textbf{Nota:}\\
$u_0 $ non è $y_0$, per esempio $\pi$ possiamo prendere solo $n$ cifre dopo la virgola, non il valore esatto. La differenza può tendere a $0$ se magari ho una procedura che approssima $\pi$.
 \begin{dimo}
	 Pensiamo $\{y_n\}$ come perturbazione di $\{u_n\}$ 
	 \[
		 u_{n+1} = u_n + h\Phi(t_n,u_n)
	 .\] 
	 \[
		 y_{n+1} = y_n + h\Phi(t_n,y_n) + h \tau_{n+1}(h)
	 .\] 
	 dove $\tay_{n+1}(h)$ è l'errore locale di troncamento.\\
	 Se chiamiamo\ \ $h[\Phi + \delta_{n+1}]$\\
	 $\delta_0 = y_0 - u_0$\\
	 $\delta_{n+1} = \tau_{n+1}(h)$ \\
	 Ripercorrendo la dimostrazione che utilizza il lemma di Gronwall discreto per la zero-stabilità, segue immediatamente
	 \[
		 |y_n - u_n| \leq (|y_0 - u_0| + nh\tau(h))e^{nh\Lambda} \ \ \ nh = t_n \leq T \ \ 0\leq n \leq N
	 .\] 
	 dove $\displaystyle\tau(h) = \sup_{n}|\tau_n(h)|$ \\
	 $ \Rightarrow  |y_n - u_n|\leq (|y_0 - u_0| + T\tau (h))e^{T\Lambda}  = \varepsilon C$ \\
	 dove $\varepsilon = |y_0-u_0 | + T\tau(h)$ e $C = e^{T\Lambda}$ e abbiamo $|y_0-u_0|$ che è l'errore sui dati, $T\tau(h)$ è la consistenza e $e^{T\Lambda}$ è la zero-stabilità\\
	 Usando infine le ipotesi\\
	 se  $| y_0 - u_0| \xrightarrow{h \rightarrow 0} 0$ e $\tau(h) \xrightarrow{ h \rightarrow 0} 0$  allora $|y_n-u_n| \xrightarrow{h \rightarrow 0} 0$ $0 \leq n\leq N$\\
	 Se inoltre $|y_0 - u_0| = O(h^s)$ e $\tau(h) = O(h^q)$ allora l'ordine di convergenza è $O(h^{\min\{q,s\}})$ (Si capisce facilemnte dalla disequaizone sopra, la somma si "mangia" l'ordine maggiore)
\end{dimo}
\textbf{Nota:}\\
Da notare come la richiesta sulla soluzione è solamente la sua esistenza.\\
Consistenza + Zero-Stabilità $ \Rightarrow $ Convergenza\\
\textbf{Analisi esplicita per il metodo di Eulero in avanti}\\
Inizialmente assumiamo $u_0 = y_0$\\
\[
	u_{n+1} = u_n + hf(t_n,u_n)
.\] 
La soluzione esatta risolve $y'(t_n) = f(t_n,y(t_n))$ scriviamo in modo più compatto  $y'_n = f(t_n,y_n)$\\
e  soddisfa
\[
	y_{n+1} = y_n  + hf(t_n,y_n) + h\tau_{n+1}(h)
.\] 
dato che c'è l'errore di troncamento\\
Definiamo 
\[
	u_{n+1}^* = y_n + hf(t_n,y_n) = y_n + h y'(t_n)
.\] 
TODO AGGIUNGI FOTO 20 marzo 11:33\\
L'errore 
 \[
	 e_{n+1} := y_{n+1} - u_{n+1} = (y_{n+1}-u_{n+1}^*) + (u_{n+1}^*- u_{n+1})
\] \[
= h\tau_{n+1}(h) + \alpha
.\] 
dove $\alpha$ è l'accumulazione degli errori precedenti.
\[
	u_{n+1}^*-u_{n+1}= y_n-u_n + h(f(t_n,y_n) - f(t_n,u_n)) = e_n + h(f(t_n,y_n) - f(t_n,u_n))
.\] 
\[
	\Rightarrow  e_{n+1} = h\tau_{n+1}(h) + e_n + h(f(t_n,y_n) - f(t_n,u_n))
.\] 
\[
	|e_{n+1}|\leq h | \tau_{n+1}(h)| + |e_n| + hL |y_n - u_n|\ \ \substack{\leq\\ \sup_n|\tau_{n+1}} \ \ h|\tau(h) | + (1 + hL)|e_n|
.\] 
$e_0 = y_0 - u_0= 0$\\
$|e_1|\leq h|\tau(h)|$, $|e_2|\leq h|\tau(h)| + (1 + hL)|e_1|\leq h|\tau(h) + (1+hL)h|\tau(h)| = h|\tau(h)|(1 + (1 + hL))$\\
$|e_n|\leq h|\tau(h)|[1 + (1 + hL) + (1 + hL)^2 + \ldots + (1 + hL)^{n-1}]$ 
 $\displaystyle = h |\tau(h)|\sum^{n-1}_{k=0}(1+hL)^k$\\
 $ \displaystyle = h|\tau(h)|\frac{1 - (1 + hL)^n}{1 - (1+hL)}= h|\tau(h)|\frac{1 - (1 + hL)^n}{-hL} =h|\tau(h)| \frac{(1 - hL)^n - 1}{hL}$\\
 e utilizzando la stima  $h\leq \frac Tn $\\
 $\displaystyle \leq \frac{(1 + \frac T n L )^n-1}{hL} \leq \frac{e^{LT}-1}{hL}$\\
 \[|e_n|\leq \cancel{h}|\tau(h)|\frac{e^{Lt}-1}{\cancel{h}L}.\]

 $h\tau_{n+1} = y_{n+1} - y_n - hf(t_n,y_n)$\\
 $h\tay_{n+1} = y_n + h y_n' + \frac {h^2}2 y''(\xi_n)- y_n-h f(t_n,y_n) = \frac {h^2}2 y''(\xi_n) \ \ t_n\leq\xi_n\leq t_{n+1}$\\
 $\tau_{n+1} = \frac h 2 y''(\xi_n)$\\
 $\displaystyle|\tau_{n+1}|\leq \frac h 2 \sup_{t}|y''(t)| = h \frac M 2$ \\
 con $M = \sup_t|y''(t)|$\\
 In conclusione
  \[
	  |e_n|\leq h\frac M 2 \frac{e^{LT} - 1}L = h C
 .\] 
 con $C = \frac M 2 \frac{e^{LT} - 1}L$ \\
 Quindi Eulero fa schifo perché è di ordine 1 ma $C$ può essere molto grande.\\
 Vedremo il caso in cui i dati non son esatti.


\end{document}
